% ── Περίληψη (Greek abstract, NTUA edition) ──
% Ακριβής απόδοση της αγγλικής Abstract.

\greekchapter{Περίληψη}

Η αυξανόμενη υιοθέτηση συσκευών του Διαδικτύου των Ιατρικών Πραγμάτων
(Internet of Medical Things, \gls{iomt}) στα περιβάλλοντα υγειονομικής περίθαλψης
έχει εισαγάγει σημαντικές προκλήσεις κυβερνοασφάλειας. Οι συσκευές αυτές---από
φορετούς αισθητήρες έως έξυπνες αντλίες έγχυσης---μεταδίδουν ευαίσθητα δεδομένα
ασθενών μέσω δικτύων που αποτελούν ολοένα και συχνότερα στόχο αντιπάλων.
Παρόλο που τα συστήματα ανίχνευσης εισβολών (Intrusion Detection Systems, \gls{ids})
που βασίζονται στη μηχανική μάθηση (Machine Learning, \gls{ml}) έχουν δείξει
δυνατότητες ανίχνευσης δικτυακών εισβολών, η αποτελεσματικότητά τους περιορίζεται
θεμελιωδώς από την ποιότητα και την αντιπροσωπευτικότητα των δεδομένων εκπαίδευσης.

Κρίσιμος περιορισμός στην τρέχουσα έρευνα είναι η εξάρτηση από μεμονωμένα,
απομονωμένα σύνολα δεδομένων που αποτυπώνουν μόνο ένα στενό τμήμα του τοπίου
επιθέσεων. Μοντέλα εκπαιδευμένα σε τέτοια δεδομένα συχνά αποτυγχάνουν να γενικεύσουν
σε διαφορετικές τοπολογίες δικτύου, διαμορφώσεις συσκευών και διανύσματα επίθεσης.
Η παρούσα διπλωματική εργασία αντιμετωπίζει αυτό το κενό προτείνοντας και υλοποιώντας
μια συστηματική μεθοδολογία για την ενοποίηση πολλαπλών ετερογενών συνόλων δεδομένων
δικτυακής κίνησης \gls{iomt} σε ένα ενιαίο, ενοποιημένο σύνολο δεδομένων, έτοιμο για
μηχανική μάθηση και κατάλληλο για την εκπαίδευση εύρωστων μοντέλων \gls{ids}.

Η μεθοδολογία επικεντρώνεται στη χρήση του MIOTTA-NPT, ενός εξειδικευμένου
εργαλείου προεπεξεργασίας δικτύου που εξάγει ένα τυποποιημένο σύνολο στατιστικών
χαρακτηριστικών από ακατέργαστες καταγραφές πακέτων, με βάση παραθυρική συνάθροιση
ρυθμιζόμενου παραθύρου. Με την επεξεργασία διαφορετικών συνόλων δεδομένων μέσω
μιας κοινής διοχέτευσης εξαγωγής χαρακτηριστικών, η θεμελιώδης πρόκληση της
ετερογένειας χαρακτηριστικών μεταξύ συνόλων δεδομένων επιλύεται στην πηγή.
Η διοχέτευση ενοποίησης ενσωματώνει επιπλέον εναρμόνιση ετικετών σε τρία επίπεδα
λεπτομέρειας, διαδικασίες διασφάλισης ποιότητας δεδομένων, κλιμάκωση χαρακτηριστικών
και ιχνηλάτηση προέλευσης του συνόλου δεδομένων, ώστε να καθίσταται δυνατή η αυστηρή
αξιολόγηση μεταξύ τομέων.

Το προκύπτον ενοποιημένο σύνολο δεδομένων διατηρεί την ποικιλομορφία των
συστατικών πηγών---συμπεριλαμβανομένων διαφορετικών τύπων επιθέσεων, πρωτοκόλλων
δικτύου και συμπεριφορών συσκευών---εξασφαλίζοντας ταυτόχρονα τη δομική συνέπεια
που απαιτείται για αποτελεσματική εκπαίδευση μοντέλων. Η έκδοση Phase~1
(8.011.534 εγγραφές παραθύρων με εναρμονισμένες ετικέτες και κλιμακωμένα
χαρακτηριστικά) είναι δημόσια διαθέσιμη στο αποθετήριο ερευνητικών δεδομένων CORA.
Τρεις ταξινομητές βασισμένοι σε δέντρα---Decision Tree, Random Forest και
XGBoost---εκπαιδεύτηκαν τόσο σε καθεστώτα εκπαίδευσης μίας πηγής όσο και σε
ενοποιημένο καθεστώς, και αξιολογήθηκαν σε διατηρημένα τμήματα δοκιμής καθώς και
στο HospNet26, μια καταγραφή κίνησης πραγματικού νοσοκομειακού δικτύου, προκειμένου
να εκτιμηθεί η γενίκευση μεταξύ τομέων. Τα πειραματικά αποτελέσματα δείχνουν ότι
τα μοντέλα που εκπαιδεύτηκαν στο ενοποιημένο σύνολο επιτυγχάνουν macro-F1
$\geq 0{,}948$ και στα δύο πεδία προέλευσης, υπερτερώντας καθοριστικά των βασικών
μοντέλων μίας πηγής (macro-F1 $\leq 0{,}74$). Η εφαρμογή στο HospNet26 αποκαλύπτει
σημαντική μετατόπιση συνμεταβλητών (covariate shift) μεταξύ εργαστηρίου και
κλινικής---το $90$--$96\%$ της καλόβουλης κλινικής κίνησης ταξινομείται εσφαλμένα
ως επίθεση, κυρίως ως DoS---παρέχοντας την πρώτη ποσοτική μέτρηση αυτού του χάσματος
και τεκμηριώνοντας συγκεκριμένες στρατηγικές επαναβαθμονόμησης για την ασφαλή
λειτουργική ανάπτυξη συστημάτων ανίχνευσης εισβολών σε περιβάλλοντα \gls{iomt}.

\bigskip
\noindent\textbf{Λέξεις Κλειδιά:} Διαδίκτυο των Ιατρικών Πραγμάτων, Σύστημα Ανίχνευσης Εισβολών, Μηχανική Μάθηση, Ενοποίηση Συνόλων Δεδομένων, Ανάλυση Δικτυακής Κίνησης, Εξαγωγή Χαρακτηριστικών, Κυβερνοασφάλεια.
